\chapternotnumbered{Conclusion}
\label{chap:Conclusion}

The primary goal of this thesis was to investigate the viability of Smoothed Particle Hydrodynamics (SPH) and the Semi-Implicit Lagrangian Voronoi Approximation method (SILVA) for simulating orographic flows. As outlined in \Cref{chap:FLOWS}, topography-driven flows are responsible for common but serious atmospheric phenomena, such as internal gravity waves, downslope windstorms, and clear-air turbulence. However, their complex dynamics make them difficult to model numerically. In \Cref{chap:MODELLING}, we first formulated the proposed problem as an initial-boundary value problem for the isentropic Euler system with ideal gas thermodynamics, and established the mathematical requirements for robust and correct modeling of these flows. We highlighted that any successful numerical method must strictly respect the stratified atmospheric background --- must be \emph{well-balanced}.

\Cref{chap:SPH} presented the theoretical background of the central method and the main contributions of this work. After revisiting the classical, Hamiltonian, and Lagrangian derivations of SPH, we proposed density-independent formulations built on \emph{potential temperature} and on the \emph{Exner function} --- to our knowledge, the first density-independent SPH cast in the thermodynamic variables native to atmospheric science. Building on these, we constructed a \emph{novel well-balanced discretization of the pressure-potential-temperature formulation} that preserves the discrete hydrostatic balance \emph{exactly, to machine precision}. \Cref{chap:IMPLEMENTATION} documented the implementation of these techniques and provides references to the publicly available code.

\Cref{chap:SILVA} was devoted to a brief overview of the SILVA method. Crucially, we have discussed how, in the current stage of development, the standard SILVA possesses serious limitations for simulating compressible flows. In \Cref{sec:SILVA_FAILURE} we showed that the exact hydrostatic balance is \emph{not preserved} and, interestingly, that the total hydrostatic residual $\abs{\grad p/ \rho + \vb{g}}$ is dominated by the boundary residual. Moreover, we have observed that while in the domain's interior the preservation can be improved by refining the mesh (but still \emph{not completely resolved}), the boundary hydrostatic residual seems unaffected by the refinement. Consequently, we have decided to not further develop SILVA.

\Cref{chap:EXPERIMENTS} put the methods to the test. Where the standard SPH formulation develops spurious velocities that reach roughly half the speed of sound within seconds and destroy the stratified background, the proposed well-balanced formulation held the hydrostatic state at machine precision and captured the emergence and early upward propagation of mountain waves. The simulations were ultimately compromised by the onset of \emph{tensile instability}, a known SPH artifact.

Summarizing, this thesis demonstrates that SPH possesses potential for atmospheric modeling, but standard formulations are insufficient. Making SPH a reliable tool for atmospheric dynamics will require the development of advanced specialized techniques. Based on our findings, future work must focus on constructing a well-balanced, density-independent $\delta$-SPH formulation to suppress tensile instability without sacrificing the exact hydrostatic balance. The use of SILVA for atmospheric modeling is, at present, limited; the natural first step would be the development of a well-balanced SILVA formulation.
